\documentclass[11pt, letterpaper]{article}

\usepackage[top=0.55in, bottom=0.55in, left=0.65in, right=0.65in]{geometry}
\usepackage{setspace}
\usepackage{parskip}
\usepackage{enumitem}
\usepackage{titlesec}
\usepackage{hyperref}
\usepackage{microtype}
\usepackage{lmodern}
\usepackage[T1]{fontenc}

\hypersetup{colorlinks=true, linkcolor=blue, urlcolor=blue}

\titleformat{\section}{\normalsize\bfseries}{}{0em}{}
\titlespacing*{\section}{0pt}{0.5em}{0.2em}

\setstretch{1.0}
\setlength{\parskip}{0.28em}
\setlength{\parindent}{0pt}

\begin{document}

%% =====================================================================
%% PAGE 1 — SPECIFIC AIMS
%% =====================================================================

\begin{center}
    {\large\bfseries Specific Aims}\\[0.15em]
    {\small Shrinithi Natarajan}
\end{center}

\vspace{0.15em}
\hrule
\vspace{0.35em}

Drug resistance causes more than 90\% of cancer deaths in patients receiving targeted therapy [1],
yet no computational framework predicts resistance by reasoning explicitly over the causal hierarchy
that governs how conflicting multi-omics signals resolve in a living cell.
Matched genomic, transcriptomic, and pharmacological profiles for thousands of cancer cell
line--drug pairs are publicly available through GDSC2 and CCLE/DepMap [2].
Deep learning baselines (DRN-CDR [3], MSDRP [4]) reconcile conflicting data layers through
learned weights that are opaque and clinically uninterpretable.
LLM-based multi-agent systems introduce natural-language reasoning but suffer from context
overwhelm and sycophantic convergence when conflicts are settled by unstructured negotiation
rather than explicit biological rules [5].
The critical gap is the absence of a framework that encodes the causal ordering of evidence
layers---DNA structural state, transcriptional activity, pathway topology, pharmacological
history---as executable, auditable logic governing inter-agent conflict resolution.
Without it, the field cannot build oncology AI whose predictions a clinician can inspect and
contest, foreclosing the leap from accuracy-optimised black boxes to accountable clinical
decision support.

The \textbf{long-term goal} is to establish structured multi-modal argumentation as the
standard paradigm for biologically faithful oncology AI. The \textbf{objective} of this
proposal is to build, evaluate, and ablate a decentralized multi-agent consensus framework
for cancer drug resistance prediction in which every prediction is governed by a formalised
Biological Hierarchy of Truth and accompanied by a full debate trace.
\textit{We hypothesise that a four-agent system (Genomics, Transcriptomics, Pharmacology,
Pathway) resolving inter-agent conflicts through an explicit five-tier axiom hierarchy will
predict drug resistance more accurately and produce more biologically faithful reasoning
traces than monolithic LLM baselines and naive majority-vote multi-agent systems.}
This hypothesis is grounded in (1) published tumour board practice, in which DNA-level
findings hierarchically override expression signals [6], and (2) a preliminary implementation
with four MCP-connected agents, a ConflictDetector, an AxiomResolver, and a DebateEngine
validated against 40 GDSC2 cases with 143 passing tests, confirming technical feasibility.
Completion will enable, for the first time, evaluation of oncology AI on reasoning
quality---hallucination rate and biological faithfulness---metrics absent from all prior work.
All four MCP servers [7,8], the LLM-agnostic client (Gemini Flash, Qwen-32B, Gemma-9B,
Llama-3.3), and the full evaluation pipeline are implemented and tested.

\textbf{Specific Aim 1: Determine whether axiom-governed multi-agent debate predicts cancer
drug resistance more accurately than monolithic LLM and classical ML baselines.}\\
\textit{Working hypothesis:} Explicit causal tier resolution prevents evidence-layer conflicts
from collapsing into statistical averaging, improving AUROC and AUPRC over a single-LLM and
a Random Forest baseline trained on the same GDSC2 features.\\
\textit{Approach:} Run the full framework and both baselines on 40 stratified GDSC2 development
cases across four free LLM backends; validate on held-out CTRPv2 cases.\\
\textit{Milestones/metrics:} Framework AUROC $\geq 0.75$ and AUPRC $\geq 0.70$; framework
exceeds both baselines on both metrics; CTRPv2 AUROC within 0.05 of development-set AUROC.

\textbf{Specific Aim 2: Determine which architectural components independently contribute
to predictive accuracy and reasoning quality through systematic ablation.}\\
\textit{Working hypothesis:} Removing the axiom hierarchy, collapsing to a single agent, or
replacing MCP-validated tool returns with free-text summaries will each independently degrade
performance, confirming the three core innovations are non-redundant.\\
\textit{Approach:} Execute six pre-specified ablations (no axioms, single agent, free-text
tools, randomised tier order, no debate rounds, majority-vote resolution) on the same 40 cases.\\
\textit{Milestones/metrics:} Each ablation shows significant degradation on $\geq 1$ of:
AUROC, hallucination rate ($<$5\% target for full framework), or biological faithfulness
score (axiom invocations consistent with established molecular biology $\geq$90\%).

The proposed work is innovative because it is the first to formalise inter-modality conflict
resolution in oncology AI as executable axiom logic [5], the first to deploy MCP as a
structural anti-hallucination layer [7,8], and the first to introduce hallucination rate and
biological faithfulness as primary evaluation metrics alongside predictive accuracy.
We expect the framework to exceed all baselines on AUROC and AUPRC, achieve a hallucination
rate below 5\%, and for each ablation to independently degrade performance.
These results will establish an open-source, model-agnostic framework that clinical teams can
deploy to generate auditable drug sensitivity predictions---advancing the NCI mission to reduce
cancer mortality through accountable AI-assisted precision oncology.

%% =====================================================================
%% PAGE 2 — REFLECTION, FEEDBACK, REFERENCES
%% =====================================================================

\newpage

\begin{center}
    {\large\bfseries Reflection, Feedback \& References}\\[0.15em]
    {\small Shrinithi Natarajan}
\end{center}

\vspace{0.15em}
\hrule
\vspace{0.35em}

\section*{Reflection}

The central structural lesson from this exercise was that the Specific Aims page functions as
a single continuous argument---gap $\to$ hypothesis $\to$ aims $\to$ payoff---and every
sentence must advance that chain or be cut. Drafting paragraph one forced precision: ``no
framework encodes causal tier ordering as executable logic'' is a falsifiable, bounded absence;
``tools lack interpretability'' is not. The discipline of writing the gap narrowly made the
objective and hypothesis easier to draft, because the scope was already constrained. The
hardest sentence was the italicised hypothesis, which had to be simple enough to read in three
seconds and specific enough to constrain the ablation design. Anchoring it to the three-way
comparison (axiom-governed vs.\ monolithic vs.\ majority-vote) gave it the necessary
testability. The rubric requirement that Aims be conceptual rather than procedural was the
most productive constraint: rephrasing from ``Run the DebateEngine on 40 cases'' to
``Determine whether axiom-governed debate outperforms baselines'' forced the aims to name
knowledge outcomes rather than actions, which clarified what the project actually claims to
show. Fitting everything to one page required collapsing qualifications and rationale into
single clauses, which improved density without losing content. The payoff paragraph benefited
most from the Wrenshall framing: leading with ``The proposed work is innovative because''
forces the innovation claim to be stated rather than implied, and ``We expect that'' rather
than ``We hope'' commits the proposal to specific, falsifiable predictions.

\section*{Feedback}

Feedback has not yet been formally solicited for this draft. The planned approach is to share
the document with my thesis supervisor and one peer reviewer with a computational biology
background, asking both to apply the rubric criteria explicitly: is the gap statement a
falsifiable absence? Does the hypothesis sentence stand alone when isolated? Are the Aims
conceptual rather than procedural? Is each expected outcome stated with a strong verb and a
measurable threshold? Reviewer comments and resulting changes will be recorded here once
collected. Two areas where feedback is anticipated: (1) whether the two Aims are sufficiently
independent---Aim 1 establishes accuracy and Aim 2 establishes component contribution, making
them logically sequential rather than parallel, which may prompt restructuring; and (2)
whether ``biological faithfulness score'' is operationally defined clearly enough for an
external reviewer unfamiliar with the axiom system to evaluate it. If flagged, a one-sentence
operational definition will be added to the Aim 2 approach sub-field.

\vspace{0.5em}
\hrule
\vspace{0.4em}

\section*{References}

{\small
\begin{enumerate}[label={[\arabic*]},leftmargin=*,itemsep=1pt,parsep=0pt]

\item Liu X et al. Multi-omics advances in understanding cancer drug resistance.
      \textit{Biomed Technol}. 2025. doi:10.1016/j.bmt.2025.100119

\item Barretina J et al. The Cancer Cell Line Encyclopedia enables predictive modelling of
      anticancer drug sensitivity. \textit{Nature}. 2012;483:603--607.

\item Saranya KR, Vimina ER. DRN-CDR: cancer drug response prediction using multi-omics.
      \textit{Comput Biol Chem}. 2024. doi:10.1016/j.compbiolchem.2024.108175

\item Zhao H et al. MSDRP: deep learning model for drug response prediction using
      multi-omics. \textit{Bioinformatics}. 2023.

\item Liu Q et al. EvoMDT: self-evolving multi-agent system for clinical decision-making.
      \textit{npj Digit Med}. 2026. doi:10.1038/s41746-025-02304-8

\item Han X et al. Multi-Agent Medical Decision Consensus Matrix for Oncology MDT.
      \textit{arXiv}:2512.14321. 2024.

\item Flotho M et al. MCPmed: MCP-enabled bioinformatics for LLM-driven discovery.
      \textit{Brief Bioinform}. 2026. doi:10.1093/bib/bbag076

\item Widjaja F et al. BioinfoMCP: MCP Interfaces in Agentic Bioinformatics.
      \textit{arXiv}:2510.02139. 2025.

\end{enumerate}
}

\end{document}
